\section{Problemas}

\question{
  Denotar a matriz $A = (a_{ij})$ do tipo $2 \times 2$ sabendo que
  $a_{ij} = i + j$.
}
\answer{
  \begin{equation}
  \begin{aligned}
  A &=
  \begin{bmatrix}
  a_{11} & a_{12} \\
  a_{21} & a_{22}
  \end{bmatrix}
  \\
  &=
  \begin{bmatrix}
  1+1 & 1+2 \\
  2+1 & 2+2
  \end{bmatrix}
  \\
  &=
  \begin{bmatrix}
  2 & 3 \\
  3 & 4
  \end{bmatrix}
  \end{aligned}
  \end{equation}
}

\question{
  Escreva a matriz $B = (b_{ij})$ do tipo $2 \times 4$, sendo que
  $b_{ij} = 3i - j$ se $i = j$ e $b_{ij} = i - 2j$ se $i \neq j$.
}
\answer{
  \begin{equation}
  \begin{aligned}
  A &=
  \begin{bmatrix}
  a_{11} & a_{12} \\
  a_{21} & a_{22} \\
  a_{31} & a_{32} \\
  a_{41} & a_{42}
  \end{bmatrix}
  \\
  &=
  \begin{bmatrix}
  3 \times 1 - 1 & 1 - 2 \times 2 \\
  2 - 2 \times 1 & 3 \times 2 - 2 \\
  3 - 2 \times 1 & 3 - 2 \times 2 \\
  4 - 2 \times 1 & 4 - 2 \times 2 
  \end{bmatrix}
  \\
  &=
  \begin{bmatrix}
  2 & -3 \\
  0 & 4  \\
  1 & -1 \\
  2 & 0  
  \end{bmatrix}
  \end{aligned}
  \end{equation}
}

\question{
  Escreva a matriz linha $D = (d_{ij})$ do tipo $1 \times 4$, tal que
  $d_{ij} = i + 4j$.
}
\answer{
  \begin{equation}
  \begin{aligned}
  D &=
  \begin{bmatrix}
  a_{11} \\
  a_{21} \\
  a_{31} \\
  a_{41} 
  \end{bmatrix}
  \\
  &=
  \begin{bmatrix}
  1 + 4 \times 1 \\
  2 + 4 \times 1 \\
  3 + 4 \times 1 \\
  4 + 4 \times 1 
  \end{bmatrix}
  \\
  &=
  \begin{bmatrix}
  5 \\
  6 \\
  7 \\
  8 
  \end{bmatrix}
  \end{aligned}
  \end{equation}
}

\question{
  Dadas as matrizes
  \[
  A =
  \begin{bmatrix}
  2 & 3  & 8 \\
  4 & -1 & 6
  \end{bmatrix}
  \quad \text{e} \quad
  B =
  \begin{bmatrix}
  5 & -7 & -9 \\
  0 & 4  & 1
  \end{bmatrix}.
  \]
  Determine $A + B$.
}
\answer{
  \begin{equation}
  \begin{aligned}
  A + B &=
  \begin{bmatrix}
    2 + 5 & 3 + (-7) & 8 + (-9) \\
    4 + 0 & -1 + 4 & 6 + 1
  \end{bmatrix}
  \\
  &=
  \begin{bmatrix}
    7 & -4 & -1 \\
    4 & 3 & 7
  \end{bmatrix}
  \end{aligned}
  \end{equation}
}

\question{
  Dadas as matrizes
  \[
  B =
  \begin{bmatrix}
  5 & -7 & -9 \\
  0 & 4  & 1
  \end{bmatrix}
  \quad \text{e} \quad
  C =
  \begin{bmatrix}
  -1 \\
  2  \\
  1
  \end{bmatrix}.
  \]
  Determine $B * C$.
}
\answer{
  \begin{equation}
  \begin{aligned}
  B*C &=
  \begin{bmatrix}
    5 & -7 & -9 \\
    0 & 4  & 1
  \end{bmatrix}
  \begin{bmatrix}
    -1 \\
    2 \\
    1
  \end{bmatrix}
  \\
  &=
  \begin{bmatrix}
    5(-1) + (-7)(2) + (-9)(1) \\
    0(-1) + 4(2) + 1(1)
  \end{bmatrix}
  \\
  &=
  \begin{bmatrix}
    -28 \\
    9
  \end{bmatrix}
  \end{aligned}
  \end{equation}
}

\question{
  Dadas as matrizes
  \[
  A =
  \begin{bmatrix}
  2 & 3  & 8 \\
  4 & -1 & 6
  \end{bmatrix}
  \quad \text{e} \quad
  B =
  \begin{bmatrix}
  5 & -7 & -9 \\
  0 & 4  & 1
  \end{bmatrix}.
  \]
  Determine $B - A$.
}
\answer{
  \begin{equation}
  \begin{aligned}
  B - A &=
  \begin{bmatrix}
    5 - 2 & -7 - 3 & -9 - 8 \\
    0 - 4 & 4 - (-1) & 1 - 6
  \end{bmatrix}
  \\
  &=
  \begin{bmatrix}
    3 & -10 & -17 \\
    -4 & 5 & -5
  \end{bmatrix}
  \end{aligned}
  \end{equation}
}

\question{
  Dada a matriz
  \[
  A =
  \begin{bmatrix}
  2 & 3  & 8 \\
  4 & -1 & 6
  \end{bmatrix}.
  \]
  Determine $-2A$.
}
\answer{
  \begin{equation}
  \begin{aligned}
  -2A &=
  -2
  \begin{bmatrix}
    2 & 3 & 8 \\
    4 & -1 & 6
  \end{bmatrix}
  \\
  &=
  \begin{bmatrix}
    -4 & -6 & -16 \\
    -8 & 2 & -12
  \end{bmatrix}
  \end{aligned}
  \end{equation}
}

\newpage
\question{
  Dada a matriz
  \[
  C =
  \begin{bmatrix}
  -1 \\
  2  \\
  1
  \end{bmatrix}.
  \]
  Determine $-C$.
}
\answer{
  \begin{equation}
  \begin{aligned}
  -C &=
  -1
  \begin{bmatrix}
    -1 \\
    2 \\
    1
  \end{bmatrix}
  \\
  &=
  \begin{bmatrix}
    1 \\
    -2 \\
    -1
  \end{bmatrix}
  \end{aligned}
  \end{equation}
}

\question{
  Calcule os valores de $m$ e $n$ para que as matrizes
  \[
  A =
  \begin{bmatrix}
  8      & 15n \\
  12 + m & 3 
  \end{bmatrix}
  \quad \text{e} \quad
  B =
  \begin{bmatrix}
  8 & 75  \\
  6 & 3 
  \end{bmatrix}
  \]
  sejam iguais.
}
\answer{
  \begin{equation}
  \begin{aligned}
  A = B &\Rightarrow
  \begin{cases}
    15n = 75 \\
    12 + m = 6
  \end{cases}
  \\
  &\Rightarrow
  \begin{cases}
    n = 5 \\
    m = -6
  \end{cases}
  \end{aligned}
  \end{equation}
}

\question{
  Calcule os valores de $x$ para que as matrizes
  \[
  A =
  \begin{bmatrix}
  7 & 8 \\
  4 & x^2
  \end{bmatrix}
  \quad \text{e} \quad
  B =
  \begin{bmatrix}
  7 & 8 \\
  4 & 10x - 25
  \end{bmatrix}
  \]
  sejam iguais.
}
\answer{
  \begin{equation}
  \begin{aligned}
  A = B &\Rightarrow x^2 = 10x - 25
  \\
  &\Rightarrow x^2 - 10x + 25 = 0
  \\
  &\Rightarrow (x - 5)^2 = 0
  \\
  &\Rightarrow x = 5
  \end{aligned}
  \end{equation}
}

\question{
  Dadas as matrizes
  \[
  A =
  \begin{bmatrix}
  1  & 2  & 3 \\
  -2 & -5 & 7 \\
  3  & 9  & -8
  \end{bmatrix}
  \quad \text{e} \quad
  X =
  \begin{bmatrix}
  x_1 \\
  x_2 \\
  x_3
  \end{bmatrix},
  \]
  efetue a multiplicação das matrizes $A$ e $X$.
}
\answer{
  \begin{equation}
  \begin{aligned}
  AX &=
  \begin{bmatrix}
    1 & 2 & 3 \\
    -2 & -5 & 7 \\
    3 & 9 & -8
  \end{bmatrix}
  \begin{bmatrix}
    x_1 \\
    x_2 \\
    x_3
  \end{bmatrix}
  \\
  &=
  \begin{bmatrix}
    x_1 + 2x_2 + 3x_3 \\
    -2x_1 - 5x_2 + 7x_3 \\
    3x_1 + 9x_2 - 8x_3
  \end{bmatrix}
  \end{aligned}
  \end{equation}
}

\question{
  Dadas as matrizes
  \[
  A =
  \begin{bmatrix}
  5  & 0  & 6 \\
  -8 & 0  & 3 \\
  -2 & 2  & 7 \\
  1  & -1 & -5
  \end{bmatrix},
  \quad
  B =
  \begin{bmatrix}
  1 & -3 & -2 & 4  \\
  7 & 8  & 5  & 9  \\
  0 & 6  & 3  & -8
  \end{bmatrix},
  \]
  determine $(AB)^T$.
}
\answer{
  \begin{equation}
  \begin{aligned}
  AB &=
  \begin{bmatrix}
    5 & 0 & 6 \\
    -8 & 0 & 3 \\
    -2 & 2 & 7 \\
    1 & -1 & -5
  \end{bmatrix}
  \begin{bmatrix}
    1 & -3 & -2 & 4 \\
    7 & 8 & 5 & 9 \\
    0 & 6 & 3 & -8
  \end{bmatrix}
  \\
  &=
  \begin{bmatrix}
    5 & -15+36 & -10+18 & 20-48 \\
    -8 & 24+18 & 16+9 & -32-24 \\
    -2+14 & 6+16+42 & 4+10+21 & -8+18-56 \\
    1-7 & -3-8-30 & -2-5-15 & 4-9+40
  \end{bmatrix}
  \\
  &=
  \begin{bmatrix}
    5 & 21 & 8 & -28 \\
    -8 & 42 & 25 & -56 \\
    12 & 64 & 35 & -46 \\
    -6 & -41 & -22 & 35
  \end{bmatrix}
  \end{aligned}
  \end{equation}

  \begin{equation}
  \begin{aligned}
  (AB)^T &=
  \begin{bmatrix}
    5 & -8 & 12 & -6 \\
    21 & 42 & 64 & -41 \\
    8 & 25 & 35 & -22 \\
    -28 & -56 & -46 & 35
  \end{bmatrix}
  \end{aligned}
  \end{equation}
}

\question{
  Dada a matriz
  \[
  A =
  \begin{bmatrix}
  2 & -7 & 1 \\
  3 & 4  & 2 \\
  5 & -9 & 6
  \end{bmatrix},
  \]
  determine $A + A^T$.
}
\answer{
  \begin{equation}
  \begin{aligned}
  A^T &=
  \begin{bmatrix}
    2 & 3 & 5 \\
    -7 & 4 & -9 \\
    1 & 2 & 6
  \end{bmatrix}
  \\
  A + A^T &=
  \begin{bmatrix}
    2 & -7 & 1 \\
    3 & 4 & 2 \\
    5 & -9 & 6
  \end{bmatrix}
  +
  \begin{bmatrix}
    2 & 3 & 5 \\
    -7 & 4 & -9 \\
    1 & 2 & 6
  \end{bmatrix}
  \\
  &=
  \begin{bmatrix}
    4 & -4 & 6 \\
    -4 & 8 & -7 \\
    6 & -7 & 12
  \end{bmatrix}
  \end{aligned}
  \end{equation}
}

\newpage
\question{
  Calcular $m$ e $n$ para que a matriz $B$ seja inversa da matriz $A$.
  \[
  A =
  \begin{bmatrix}
  m  & -22 \\
  -2 & n
  \end{bmatrix}
  \quad \text{e} \quad
  B =
  \begin{bmatrix}
  5 & 22 \\
  2 & 9
  \end{bmatrix}.
  \]
}
\answer{
  \begin{equation}
  \begin{aligned}
  AB &=
  \begin{bmatrix}
    m & -22 \\
    -2 & n
  \end{bmatrix}
  \begin{bmatrix}
    5 & 22 \\
    2 & 9
  \end{bmatrix}
  =
  \begin{bmatrix}
    5m - 44 & 22m - 198 \\
    -10 + 2n & -44 + 9n
  \end{bmatrix}
  \end{aligned}
  \end{equation}

  Como $B = A^{-1}$, então $AB = I$:

  \begin{equation}
  \begin{aligned}
  \begin{bmatrix}
    5m - 44 & 22m - 198 \\
    -10 + 2n & -44 + 9n
  \end{bmatrix}
  =
  \begin{bmatrix}
    1 & 0 \\
    0 & 1
  \end{bmatrix}
  \end{aligned}
  \end{equation}

  \begin{equation}
  \begin{aligned}
  5m - 44 &= 1 \\
  9n - 44 &= 1
  \end{aligned}
  \qquad
  \Rightarrow
  \qquad
  \begin{aligned}
  m &= 9 \\
  n &= 5
  \end{aligned}
  \end{equation}
}

\question{
  Responda:
  \begin{enumerate}
    \item[(a)] Se $A$ é uma matriz simétrica, então $A - A^T$ é?
    \item[(b)] Se $A$ é uma matriz triangular superior, então $A^T$ é?
    \item[(c)] Se $A$ é uma matriz diagonal, então $A^T$ é?
  \end{enumerate}
}
\answer{
  \begin{enumerate}
    \item[(a)] Como $A$ é simétrica, então $A = A^T$. Logo:
    \[
    A - A^T = A - A = 0
    \]
    portanto, resulta em uma matriz nula.

    \item[(b)] Se $A$ é triangular superior, então $A^T$ é triangular inferior.

    \item[(c)] Se $A$ é diagonal, então $A^T$ continua sendo diagonal.
  \end{enumerate}
}
